\documentclass[11pt]{article}
\usepackage[T1]{fontenc}
\usepackage{lmodern,amsmath,amssymb,amsthm,booktabs,array}
\usepackage[margin=1in]{geometry}
\usepackage[hidelinks]{hyperref}
\newtheorem{theorem}{Theorem}
\newtheorem{lemma}[theorem]{Lemma}
\newtheorem{corollary}[theorem]{Corollary}
\newcommand{\Z}{\mathbb Z}
\newcommand{\N}{\mathbb N}
\title{Halving-partition nonregularity and\protect\\
dyadic extensions of Alekseyev (2026)\protect\\
arXiv:2611.01243v1, Section 6, paragraph 2, p.\ 18}
\author{Haobo Ma\\{\small ChronoAI Pte Ltd, Singapore}
\and Wenlin Zhang\\{\small National University of Singapore}\\
{\small ChronoAI Pte Ltd, Singapore}}
\date{10 November 2026}
\begin{document}
\maketitle
\begin{abstract}
Extending Alekseyev's (2026) fixed-length enumeration, we answer negatively
the question of whether the halving-partition count is
$2$-regular, posed in Section~6, second paragraph, page~18, of
\emph{Halving Partitions and Hyperbinary Expansions}
(arXiv:2611.01243v1, 2026).
An exact classification at $n=2^m+1$ gives a step function whose first
difference has arbitrarily long zero blocks followed by a one.
The dyadic length and total-part identities specialize his Theorem~5;
the list classification supplies an alternative proof.
A finite decimation family would impose a monic linear recurrence on this
difference, which is impossible. The argument also excludes finite
decimation families over every nonzero commutative ring and shows that the
count modulo any integer $q\ge2$ is not $2$-automatic. We also give explicit
dyadic deficit and digit-weight formulas, a flat length distribution,
and the same nonregularity obstruction in every base $2^s$, $s\ge1$.
\end{abstract}

\section{The question and the counting convention}
Alekseyev's Section~6 question \cite[p.~18]{alekseyev} concerns a finite
family of integer sequences containing $a(n)$ and closed, under integer
linear combinations with fixed coefficients, under both substitutions
$n\mapsto2n$ and $n\mapsto2n+1$. We show that such a family does not exist.
This is a question about finite generation, not about whether the sequence
has a finite set of values.

We use Definition~1 of \cite[p.~1]{alekseyev}. A halving partition of $n>0$
is a finite list $(p_1,\ldots,p_\ell)$ of positive integers with sum $n$,
where $p_0=n$ and
\[
 p_i\in\{\lfloor p_{i-1}/2\rfloor,\lceil p_{i-1}/2\rceil\}
 \quad(1\le i\le\ell).
\]
Coincident halves are one value, not two choices. These lists are
nonincreasing; repeated parts, including terminal ones, are retained.
The empty list is the unique partition of zero. Let $a_\ell(n)$ count
partitions of length $\ell$ and $a(n)=\sum_\ell a_\ell(n)$.
This is OEIS A349552 \cite{oeis}. OEIS A349553 records the least input
with a prescribed value of this count.

For a commutative ring $R$ with $1\ne0$, call an $R$-valued sequence
\emph{$2$-regular over $R$} if it belongs to a finite family whose two
decimations are fixed $R$-linear combinations of that family. Enumerating
the family gives a column $F(n)\in R^d$, with the sequence as first
coordinate, and matrices $M_0,M_1\in R^{d\times d}$ such that
\begin{equation}\label{eq:family}
 F(2n)=M_0F(n),\qquad F(2n+1)=M_1F(n)\quad(n\ge0).
\end{equation}
Conversely these identities give such a family. Neither linear
independence of its members nor invertibility of the matrices is assumed.
The integer case is the question in \cite{alekseyev}. For the standard
complex linear-representation formulation, see
\cite[Definition~1, p.~24:3]{heuberger}; here the coefficient ring is
specified explicitly.

\begin{theorem}\label{thm:main}
For every nonzero commutative ring $R$, the sequence
$a_R(n)=a(n)1_R$ is not $2$-regular over $R$.
\end{theorem}

\section{The exact dyadic profile}
Put $T_r=2^r+r+1$ for $r\ge1$ and, for $m\ge2$, put
\[
 R_m=\#\{r\ge1:T_r\le m\}.
\]
The set being counted is either empty or $\{1,\ldots,R_m\}$.

\begin{lemma}\label{lem:profile}
For $m\ge2$, the length enumerator is
\begin{equation}\label{eq:polynomial}
 \sum_{\ell\ge0}a_\ell(2^m+1)z^\ell
 =z^m(1+z+z^2)+\sum_{r=1}^{R_m}z^{m-r}.
\end{equation}
In particular,
\begin{equation}\label{eq:count}
 b_m:=a(2^m+1)=3+R_m.
\end{equation}
\end{lemma}
\begin{proof}
Starting at $2^m+1$, the first $m$ positions of a positive halving chain
follow one of two branches. As long as the upper branch is taken, position
$i$ has value $2^{m-i}+1$. The first lower branch gives $2^{m-i}$;
subsequent values are forced powers of two until reaching $1$, after
which all positive values are $1$.

More precisely, fix a length $\ell$. Its baseline word has entry
$2^{m-i}$ at positions $1\le i\le\min(m,\ell)$ and entry $1$ at positions
$i>m$. Every positive chain of this length is obtained by adding one to
the first $j$ entries, for a unique
$0\le j\le\min(m,\ell)$. To see exhaustiveness, observe that from an
odd value $2^s+1$ the two halves are $2^{s-1}$ and $2^{s-1}+1$;
once the former occurs, a positive half of a power of two is forced.
Conversely every displayed word obeys the halving rule. Distinct $j$
give different words, so this is a classification without multiplicity.

For $\ell\le m$ the baseline sum is $2^m-2^{m-\ell}$. The total-sum
condition is therefore
\begin{equation}\label{eq:short}
 j=2^{m-\ell}+1.
\end{equation}
For $\ell\ge m$ the baseline sum is $2^m-1+(\ell-m)$, so the condition is
\begin{equation}\label{eq:long}
 j+(\ell-m)=2.
\end{equation}
When $\ell<m$, set $r=m-\ell\ge1$. Equation~\eqref{eq:short} is
admissible exactly when $2^r+1\le\ell=m-r$, or $T_r\le m$;
it then gives one word of length $m-r$.
When $\ell\ge m$, equation~\eqref{eq:long} gives exactly
$(\ell,j)=(m,2),(m+1,1),(m+2,0)$, all admissible because $m\ge2$.
These disjoint cases prove \eqref{eq:polynomial} and \eqref{eq:count}.
\end{proof}

These enumerative identities also follow directly from Alekseyev's
Theorem~5 \cite[pp.~5--6]{alekseyev}. His recurrence (12), also
\cite[Proposition~2.2]{mcconville}, gives
$H_{2^j}(z)=z+\cdots+z^{j+1}$. For $1\le k<m$, his parameters are
$Q_k(2^m+1)=2^{m-k}$ and $M_k(2^m+1)=2^k$, so the coefficient is one
exactly when $k\ge2^{m-k}+1$. For $k\ge m$, it is
$[z^{k-1}]H_{2^m}(z)$, giving $k=m,m+1,m+2$.
Thus the length spectrum and its total-part sum are explicit
specializations of his prior theorem.

For example, $b_2=b_3=3$, $b_4=b_5=b_6=4$, and
$b_7=\cdots=b_{11}=5$. At $n=17$, the four lengths are $3,4,5,6$.

\begin{corollary}\label{cor:consequences}
For $m\ge2$, the total number of parts across all halving partitions of
$2^m+1$ is
\begin{equation}\label{eq:parts}
 \sum_\ell\ell\,a_\ell(2^m+1)
 =(R_m+3)m+3-\frac{R_m(R_m+1)}2.
\end{equation}
For $m\ge4$, with $h=\lfloor\log_2m\rfloor$, one has
$h+2\le b_m\le h+3$. Thus $b_m=\log_2m+O(1)$.
\end{corollary}
\begin{proof}
Differentiate \eqref{eq:polynomial} and set $z=1$ to obtain
\eqref{eq:parts}. If $T_r\le m$, then $2^r\le m$, hence $R_m\le h$.
Since $h\ge2$ and $h\le2^{h-1}$, we have
$T_{h-1}=2^{h-1}+h\le2^h\le m$, giving $R_m\ge h-1$.
\end{proof}

\begin{corollary}\label{cor:flat}
For $m\ge2$, precisely one partition of $2^m+1$ has each length in
$[m-R_m,m+2]$, and no other length occurs. Thus the full length
distribution is flat, symmetric and unimodal. For every modulus $d\ge1$,
the counts in length classes modulo $d$ differ pairwise by at most one.
For a uniformly chosen
partition its length has mean $m+1-R_m/2$ and variance
$((R_m+3)^2-1)/12$.
\end{corollary}
\begin{proof}
The thresholds $T_r$ increase strictly, so the exponents in
\eqref{eq:polynomial} are consecutive and each coefficient is one.
An interval of consecutive integers occupies each residue class modulo
$d$ either $\lfloor(R_m+3)/d\rfloor$ or $\lceil(R_m+3)/d\rceil$ times.
The mean is the midpoint of these $R_m+3$ integers. Translating them
to $0,\ldots,R_m+2$ and using the sums of integers and their squares
gives variance $(K^2-1)/12$ for $K=R_m+3$.
\end{proof}

\begin{corollary}\label{cor:deficit}
Let $s$ be Stern's sequence, with $s(1)=1$, $s(2n)=s(n)$ and
$s(2n+1)=s(n)+s(n+1)$ for $n\ge1$, and put $d(n)=s(n)-a(n)$.
For $m\ge2$ and $v\ge0$,
\[
 d\bigl(2^v(2^m+1)\bigr)=m-2-R_m.
\]
The odd part has $m-1$ binary zeros. Thus for this family with exactly
$z\ge1$ zero digits the deficit is $z-1-R_{z+1}$.
\end{corollary}
\begin{proof}
The Stern recurrence gives $s(2^m+1)=m+1$ by induction, starting with
$s(3)=2$. Subtract \eqref{eq:count}. Both $a(2n)=a(n)$
\cite[Lemma~2]{alekseyev} and $s(2n)=s(n)$ preserve the deficit.
\end{proof}

\begin{corollary}\label{cor:offset}
Let $A_\rho(n)$ count positive halving chains starting from the virtual
predecessor $n$ with total $n+\rho$, as in \cite[Section~4]{alekseyev}.
For $m\ge2$ and $\rho\ge0$,
\[
 A_\rho(2^m+1)=1+\min(m,\rho+2)
                  +\#\{r\ge1:T_r+\rho\le m\}.
\]
In particular, for $m\ge3$ these counts first remain equal to $m+1$
at offset $\rho=m-2$, agreeing with the source's exact threshold
\cite[Theorem~17, Corollary~18]{alekseyev}.
\end{corollary}
\begin{proof}
For $\ell<m$, the classified word has the required total exactly when
$j=2^{m-\ell}+1+\rho\le\ell$, giving the displayed threshold count.
For $\ell\ge m$, the condition becomes $j+(\ell-m)=\rho+2$.
Each $0\le j\le\min(m,\rho+2)$ gives exactly one length and word.
At $\rho\ge m-2$ the threshold set is empty and the long count is
$m+1$; at $\rho=m-3\ge0$ it is $m$.
\end{proof}

For the refinement in Section~6, third paragraph, of \cite{alekseyev},
set $q_i=p_{i-1}-2p_i+1$ and let $w(P)$ count the indices with $q_i=1$.
This is exactly the number of even predecessors in the chain. Define
$W_n(z,t)=\sum_P z^{\ell(P)}t^{w(P)}$, summing over halving partitions.
\begin{corollary}\label{cor:weighted}
For $m\ge2$, the sum-constrained digit refinement is
\begin{align*}
 W_{2^m+1}(z,t)={}&z^m t^{\max(0,m-3)}
       +z^{m+1}t^{m-2}+z^{m+2}t^{m-1}\\
       &+\sum_{r=1}^{R_m}z^{m-r}t^{\max(0,m-2^r-r-2)}.
\end{align*}
For $v\ge0$, $W_{2^v(2^m+1)}(z,t)=(zt)^vW_{2^m+1}(z,t)$.
\end{corollary}
\begin{proof}
In the classification of Lemma~\ref{lem:profile}, the word with length
$\ell<m$ has upper-prefix length $j=2^{m-\ell}+1$.
Its predecessors through position $\min(\ell,j+1)$ are odd; the subsequent
predecessors through position $\ell$ are even. Hence its weight exponent
is $\max(0,\ell-j-1)$. For the three long words $(\ell,j)=(m,2)$,
$(m+1,1)$ and $(m+2,0)$ the exponents are respectively
$\max(0,m-3)$, $m-2$ and $m-1$. The case $m=2$, $j=m$, stops
before any further even predecessor. These counts give the formula.
Doubling the total prepends its unique first part, the previous total;
this adds one part and one even predecessor.
\end{proof}

\section{A recurrence cannot have these gaps}
\begin{proof}[Proof of Theorem~\ref{thm:main}]
Suppose \eqref{eq:family} represents $a_R$. Repeated even decimation
and then odd decimation give, for $m\ge1$,
\[
 a_R(2^m+1)=e_1^{\mathsf T}M_1M_0^{m-1}F(1),
\]
where $e_1^{\mathsf T}$ selects the first coordinate.
Write the characteristic polynomial of $M_0$ as
$X^d+c_{d-1}X^{d-1}+\cdots+c_0$.
The Cayley--Hamilton identity, valid over a commutative ring, gives
\begin{equation}\label{eq:recurrence}
 b_{m+d}+\sum_{i=0}^{d-1}c_i b_{m+i}=0\quad(m\ge2),
\end{equation}
with the integer values mapped into $R$.
Subtracting consecutive identities shows that
$\delta_m=b_{m+1}-b_m$ obeys the same monic forward recurrence.

By \eqref{eq:count},
\[
 \delta_m=
 \begin{cases}
  1_R,&m=2^r+r\text{ for some }r\ge1,\\
  0,&\text{otherwise}
 \end{cases}
 \qquad(m\ge2).
\]
The successive positions $J_r=2^r+r$ satisfy
$J_{r+1}-J_r=2^r+1$. Thus $\delta_m$ has $2^r$ consecutive zeros
at $J_r+1,\ldots,J_{r+1}-1$, followed by $\delta_{J_{r+1}}=1_R$.
Choose $r$ with $2^r\ge d$. Starting with the first $d$ zeros,
the monic recurrence forces every subsequent term to be zero by
induction. This contradicts $1_R\ne0$. Only forward propagation is used;
there is no division by a recurrence coefficient.
\end{proof}

\begin{corollary}\label{cor:automatic}
For every integer $q\ge2$, the sequence $a(n)\bmod q$ is not
$2$-automatic.
\end{corollary}
\begin{proof}
If a finite automaton with output generates this sequence from the binary
digits of $n$, then its $2$-kernel
\[
 \{n\mapsto a(2^e n+t)\bmod q:e\ge0,\ 0\le t<2^e\}
\]
is finite. Indeed, for each binary suffix its output is a map from the
finite state set to the finite output set, so only finitely many such
maps occur; $n=0$ adds at most a separate initial value. Both decimations
send every kernel member to another kernel member. Viewed as
$\Z/q\Z$-valued sequences, these members therefore form a finite
decimation family, contrary to Theorem~\ref{thm:main}.
\end{proof}

\begin{corollary}\label{cor:bases}
For every $s\ge1$, $a_R$ has no finite base-$2^s$ decimation family over
a nonzero commutative ring $R$. For every $q\ge2$, $a(n)\bmod q$ is
not $2^s$-automatic.
\end{corollary}
\begin{proof}
Suppose a finite family $f_i$ is closed linearly under
$n\mapsto2^s n+u$, $0\le u<2^s$. Adjoin the finitely many sequences
$f_i(2^j n+t)$, $0\le j<s$, $0\le t<2^j$. Binary decimation
increases $j$ by one and replaces $t$ by $t+2^j b$, $b\in\{0,1\}$.
At $j=s-1$ the assumed base-$2^s$ identities return a linear combination
of the original family. This yields a binary family, contradicting
Theorem~\ref{thm:main}. A finite base-$2^s$ automaton gives a finite
base-$2^s$ kernel, just as in Corollary~\ref{cor:automatic}.
\end{proof}

\section{Coverage and verification}
The exact source question is Section~6, second paragraph, page~18, of
\cite{alekseyev}; its status word is ``question.'' The same section
also discusses the following distinct directions.

\smallskip
\noindent\begin{tabular}{@{}p{0.42\linewidth}p{0.53\linewidth}@{}}
\toprule
Source item & Coverage here\\
\midrule
Section 6, paragraph 2: finite integer decimation family & Answered: no;
Theorem~\ref{thm:main} with $R=\Z$.\\
Definition 1; Section 5 and Example 23: counts by length & Proved here
at $2^v(2^m+1)$: the full distribution is flat
(Corollary~\ref{cor:flat}, shifted by $v$). General unimodality is
not asserted in the source; Example 23 gives an obstruction.\\
Section 6, paragraph 1: deficit with more binary zeros &
Partial answer: explicit for $2^v(2^m+1)$ in
Corollary~\ref{cor:deficit}; the general formula is not covered.\\
Section 4, Theorem 17 and Corollary 18: eventual chain counts &
Explicit dyadic formula in Corollary~\ref{cor:offset}; its threshold
is a specialization of the prior proof by Alekseyev (2026).\\
Section 6, paragraph 1: maximum at fixed binary length &
Not determined by the dyadic classification.\\
Section 6, paragraph 3; Appendix (A14): weighted counts &
Partial answer: exact $W_n(z,t)$ on the dyadic family in
Corollary~\ref{cor:weighted}; general $n$ not covered.\\
Section 6, paragraph 3: division by larger bases &
Not covered; changing the division rule changes the partitions.\\
Theorem 5: dyadic counts and total parts & Reproved; prior proof
Alekseyev (2026), Theorem 5 by specialization. The list proof also gives
the digit-weight refinement.\\
OEIS A349553: least input with a specified count & Prior proof cited:
Alekseyev (2026), Corollary 15 gives $a(2^k-1)=k$, so the least input
exists and is at most $2^k-1$. Exact minima are not covered.\\
OEIS A349552 formulas; source Lemma 2 and Corollary 15 &
Prior proof cited: Alekseyev (2026), even dilation and Mersenne counts. These are
existing identities, not new claims here.\\
Ring, modular and base-$2^s$ extensions & Not posed in the source;
proved here in Theorem~\ref{thm:main} and
Corollaries~\ref{cor:automatic}, \ref{cor:bases}.\\
\bottomrule
\end{tabular}

\medskip
The written arguments prove all statements above, including the ring and
modular extensions. The pinned Lean declaration
\href{https://github.com/the-omega-institute/trureturing/blob/1154841a3345ca579e0d56fc4ceebee38814557c/D5/S1/Recurrence/Partitions/HalvingPartitionNonregularity.lean}
{\texttt{HalvingPartitionNonregularity.result}}
states integer nonexistence in both matrix and arbitrary finite-family
forms, their equivalence, equality of the list and composition counts
for every $n$, and $a(0)=1$. Its proof contains the dyadic classification
and count formula. The ring extension, length polynomial, total-part
formula, growth estimate, flat length distribution, deficit, digit-weight
formula, offset-count formula and base-$2^s$ conclusions are supplied by
the written proofs here.
The accompanying \texttt{verify.py} gives
finite exact checks of the counting formulas and examples; its numerical
ranges are recorded in \texttt{verification.json}.

The source's Lemma~2 and Theorems~5 and~8 already give even dilation,
hyperbinary coefficients and the sharp fixed-length bound; our dyadic
classification makes their specialized consequences explicit. No claim
of priority is made. The general deficit, fixed-digit maximum and
larger-division-base questions remain outside these formulas.

\begin{thebibliography}{9}
\bibitem{alekseyev}
M.~A. Alekseyev,
\emph{Halving Partitions and Hyperbinary Expansions},
arXiv:2611.01243v1 (2026), Definition~1, p.~1; Section~6, p.~18.
\href{https://arxiv.org/abs/2611.01243v1}{arXiv:2611.01243v1}.
\bibitem{oeis}
OEIS Foundation Inc.,
\emph{The On-Line Encyclopedia of Integer Sequences}, A349552,
``Number of halving partitions of $n$,'' introduced by C.~Kimberling
(2021), and A349553, ``Least $k$ such that $n$ is the number of halving
partitions of $k$''; consulted 10 November 2026.
\url{https://oeis.org/A349552}; \url{https://oeis.org/A349553}.
\bibitem{heuberger}
C.~Heuberger, D.~Krenn and T.~Lechner,
\emph{Analysis of Regular Sequences: Summatory Functions and
Divide-And-Conquer Recurrences},
LIPIcs \textbf{302}, AofA 2024, 24:1--24:14.
\href{https://doi.org/10.4230/LIPIcs.AofA.2024.24}
{doi:10.4230/LIPIcs.AofA.2024.24}.
\bibitem{mcconville}
T.~McConville, J.~Propp and B.~E. Sagan,
\emph{Hyperbinary partitions and $q$-deformed rationals},
arXiv:2508.20026v1 (2025), Proposition~2.2.
\href{https://arxiv.org/abs/2508.20026v1}{arXiv:2508.20026v1}.
\end{thebibliography}

\bigskip\noindent
\small\copyright\ 2026 Haobo Ma and Wenlin Zhang.
Licensed under Creative Commons Attribution 4.0 International
(CC BY 4.0).
\end{document}
